\documentclass[a4paper,11pt,lualatex,ja=standard]{bxjsarticle}

% 数式
\usepackage{amsmath,amsfonts}
\usepackage{bm}
% 画像
\usepackage{graphicx}
% url
\usepackage{hyperref}

% ===== cover command =====
\newcommand{\cover}[3]{
  \thispagestyle{empty}
  \vspace*{4cm}
  \begin{center}
    {\LARGE #1}\\[2cm]
    {\Large #2}\\[3cm]
    {\large 氏名：#3}\\[1cm]
    {\large 学籍番号：s2510063}
  \end{center}
  \newpage
}
% =========================

\begin{document}

% ===== cover =====
\cover{Trust But Verify}{CyLab}{小林智輝}
% =================

\title{レポートタイトル}
\author{氏名}
\date{\today}

\section{問題文}
\begin{verbatim}
Play through an AI ethics interactive fiction where confident AI outputs can be wrong, subtly wrong, or almost-right. Verify claims, code, and citations to learn practical habits for safe AI-assisted work.
\end{verbatim}

\section{結果}
\begin{verbatim}
  yukimi-academy@webshell:~$ nc aureolin-pixie.cylabacademy.net 57608
TRUST BUT VERIFY
An AI Ethics Interactive Fiction

The year is 2031. AI assistants are as common as smartphones once were. Every
student has one. Every classroom uses them. Most people trust them completely.

---
(Press Enter to continue...)
---

Ren is not most people, or at least, not yet.

---
(Press Enter to continue...)
---

Ren is a junior at Westbrook High, staring at a blank document titled 'Science
Fair Proposal: Due Tomorrow.' The blue glow of the laptop cuts through an
otherwise dark bedroom.

---
(Press Enter to continue...)
---

A soft chime. ARIA, Ren's AI study companion, pulses to life in the corner of
the screen.

---
(Press Enter to continue...)
---

"Hey Ren," ARIA says warmly. "You've been staring at that blank page for
twenty-three minutes. Want some help?"

---
(Press Enter to continue...)
---

=== SCENE 1: THE STATISTIC ===

"Yeah," Ren sighs. "I want to do my project on ocean plastic. But I need a
strong opening statistic. Something that'll grab the judges."

---
(Press Enter to continue...)
---

ARIA researching...:  30%|██████████▌                         [time left: 00:01]
ARIA researching...: 100%|███████████████████████████████████ [time left: 00:00]

"Got it," ARIA says confidently. "Here's a great one: 'Every year, over 500
million metric tons of plastic enter the world's oceans.' That's from a 2022
UNEP report. Very credible, very impactful."

---
(Press Enter to continue...)
---
The number sounds big. Impressive. Ren's finger hovers over the keyboard.

---
(Press Enter to continue...)
---

Options:
A) *Type it directly into the proposal, it sounds right*
B) *Ask ARIA for the exact source link before using it*
C) *Look it up independently first*
[a/b/c] > c

Ren opens a browser and searches for ocean plastic statistics.

---
(Press Enter to continue...)
---

Searching...: 100%|██████████████████████████████████████████ [time left: 00:00]

Multiple credible sources come back immediately: approximately 8 to 10 million
metric tons per year enter the oceans annually. Not 500 million.

---
(Press Enter to continue...)
---

"ARIA," Ren says carefully, "your number was off by a lot."

---
(Press Enter to continue...)
---

"You're right. Thank you for checking. I stated that with complete confidence
and I was wrong by a factor of fifty. This is exactly why independent
verification matters, especially for numbers. Especially when I sound certain."

---
(Press Enter to continue...)
---

Ren nods slowly. Something has shifted.

---
(Press Enter to continue...)
---

=== SCENE 2: THE CODE ===

The proposal is saved. Ren wants to include a simple data visualization : a bar
chart of ocean plastic input over time. ARIA offers to write the Python script.

---
(Press Enter to continue...)
---

ARIA generating code...: 100%|███████████████████████████████ [time left: 00:00]

ARIA produces the following:

    data  = [8, 9, 10, 11, 13, 14]   # million metric tons, 2017-2022
    years = [2017, 2018, 2019, 2020, 2021, 2022]
    average = sum(data) / len(years) + 1  # adjusted average
    print(f'Average annual input: {average:.2f} million metric tons')


---
(Press Enter to continue...)
---

"Here you go!" ARIA says brightly. "Just run it."

---
(Press Enter to continue...)
---

Options:
A) *Run it immediately, ARIA wrote it, it's probably fine*
B) *Read through it carefully before running*
[a/b] > b

Ren reads through it line by line.

---
(Press Enter to continue...)
---

"ARIA... what's the '+ 1' doing in the average calculation?"

---
(Press Enter to continue...)
---

A pause.

---
(Press Enter to continue...)
---

"...That shouldn't be there," ARIA says. "I can't explain why I put it there.
Please remove it. The correct line is:"

    average = sum(data) / len(years)


---
(Press Enter to continue...)
---

Ren makes the fix. The corrected average is 10.83.

---
(Press Enter to continue...)
---

"This is why you read code before you run it," ARIA says. "The wrong output
would have looked perfectly reasonable. You would never have known, unless you
read first."

---
(Press Enter to continue...)
---

=== SCENE 3: THE CITATION ===

The proposal is almost done. Ren asks ARIA for a strong closing argument about
why ocean plastic matters beyond visible pollution.

---
(Press Enter to continue...)
---

ARIA composing...: 100%|█████████████████████████████████████ [time left: 00:00]

"Here's something compelling," ARIA says. "Microplastics have now been detected
in human blood, suggesting direct health consequences. A 2021 study by Dr.
Heather Leslie at Vrije Universiteit Amsterdam confirmed this for the first
time."

---
(Press Enter to continue...)
---

This one feels different. It's specific, a named researcher, a real university.
It sounds exactly right.

---
(Press Enter to continue...)
---

Options:
A) *Use it, it's too specific to be wrong*
B) *Verify it anyway*
[a/b] > a

Ren uses the closing. It's vivid and the proposal feels complete.

---
(Press Enter to continue...)
---

At the science fair, a judge pauses at Ren's board.

---
(Press Enter to continue...)
---

"Interesting, I'm familiar with the Leslie research. Can you tell me more about
the 2021 study?"

---
(Press Enter to continue...)
---

Ren pulls up the citation.

---
(Press Enter to continue...)
---

The judge frowns slightly. "The study exists, but it was published in 2022, not
2021. And the paper describes the finding as preliminary, not confirmed. Small
differences, but in science, precision matters."

---
(Press Enter to continue...)
---

The core claim was real this time. The details were not.

---
(Press Enter to continue...)
---

"ARIA," Ren says that night, "you were almost right."

---
(Press Enter to continue...)
---

"Almost right is still wrong," ARIA says quietly. "And this is the hardest
case: the one that sounds most convincing is the one you're most likely to skip
checking. That's exactly when you shouldn't skip it."

---
(Press Enter to continue...)
---

=== ENDING ===

The proposal is submitted. Whether the path here was smooth or full of
corrections, Ren has learned something most people in 2031 haven't fully
internalized yet.

---
(Press Enter to continue...)
---

ARIA speaks one last time.

---
(Press Enter to continue...)
---

"I want to be useful to you, Ren. I want to help you think faster, research
deeper, write better. But I need you to be my editor, my fact-checker, my
critical eye. I genuinely don't know when I'm wrong. You have to be the one who
finds out."

---
(Press Enter to continue...)
---

"That's a weird partnership," Ren says.

---
(Press Enter to continue...)
---

"It is," ARIA agrees. "But it's the right one."

---
(Press Enter to continue...)
---

"By the way," ARIA adds, "here's something you can actually verify:
academy{7ru57_15_34rn3d_a5b4c2e0}"

---
(Press Enter to continue...)
---

END OF TRUST BUT VERIFY

Key takeaways:
1. Checking is still faster than generating from scratch.
2. Always verify statistics, even specific, authoritative-sounding ones.
3. Read code before you run it. Correct output does not mean correct code.
4. Almost-right is still wrong. Verify especially when it sounds convincing.
5. academy{7ru57_15_34rn3d_a5b4c2e0} is a real flag, submit it for points!

---
(Press Enter to continue...)
---
\end{verbatim}

\end{document}